% Generated draft: requires mathematical and source review.
\documentclass[11pt]{article}
\usepackage[margin=25mm]{geometry}
\usepackage{fontspec}
\setmainfont{Latin Modern Roman}
\usepackage{amsmath,amssymb,mathtools,unicode-math}
\setmathfont{Latin Modern Math}
\usepackage{enumitem,needspace,xurl,hyperref}
\setlength{\parindent}{0pt}
\setlength{\parskip}{5pt}
\setlength{\emergencystretch}{3em}
\newcommand{\N}{\mathbb{N}}
\newcommand{\Z}{\mathbb{Z}}
\newcommand{\Q}{\mathbb{Q}}
\newcommand{\R}{\mathbb{R}}
\newcommand{\C}{\mathbb{C}}
\newcommand{\F}{\mathbb{F}}
\newcommand{\A}{\mathbb{A}}
\newcommand{\PP}{\mathbb{P}}
\newcommand{\E}{\mathbb{E}}
\newcommand{\eps}{\varepsilon}
\newcommand{\dd}{\,\mathrm{d}}
\begin{document}
\section*{UnsolvedMath ID: 30005024}
\textbf{Source problem code:} catalogue entry
\subsection{Definitions and mathematical statement}
The supplied question does not specify the index graph, the invariance group, or the allowed state spaces. Consequently, the following is a parameterized formulation; it does not settle which parameters the source intends to fix or quantify over.

For a specified countable, connected, locally finite graph \(G=(V,E)\), a specified transitive group \(\Gamma\leq\operatorname{Aut}(G)\), and a specified measurable state space \(S\), let \(X=(X_v)_{v\in V}\) be an \(S\)-valued process. Write \(d_G\) for graph distance and \(B_G(v,r)=\{w\in V:d_G(v,w)\leq r\}\). The process is \emph{invariant} if its law is unchanged by every element of \(\Gamma\). It is \emph{finitely dependent} if there exists an integer \(k\geq0\) such that, for all subsets \(U,W\subseteq V\) with \(d_G(U,W)>k\), the sigma-fields \(\sigma(X_u:u\in U)\) and \(\sigma(X_w:w\in W)\) are independent.

An \emph{i.i.d. representation as a finitary factor} consists of an i.i.d. process \(Y=(Y_v)_{v\in V}\), with some permitted input state space and common probability law, and a measurable \(\Gamma\)-equivariant map \(\Phi\) such that \(\Phi(Y)\) has the law of \(X\). Finitariness means that each output coordinate is determined almost surely by the input in a finite graph ball. More precisely, on a fixed invariant set \(D\) of input configurations of probability one, a coding radius \(R_v(y)\) satisfies
\[
 y'\in D,\quad y'|_{B_G(v,R_v(y))}=y|_{B_G(v,R_v(y))}
 \quad\Longrightarrow\quad
 \Phi(y')_v=\Phi(y)_v,
\]
with \(R_v(Y)<\infty\) almost surely. The radii are understood to be measurable and chosen equivariantly. Such a representation has \emph{finite expected coding radius} if \(\mathbb E[R_o(Y)]<\infty\) at a fixed vertex \(o\); transitivity and equivariance make this condition independent of the vertex.

\textbf{Question, conditional on the intended indexing framework and allowed state spaces.} Does there exist an invariant, finitely dependent process \(X\) such that no i.i.d. input process and equivariant finitary factor map representing its law admit a coding radius of finite expectation? Equivalently, the requested property is
\[
 \text{for every permitted i.i.d. input }Y\text{ and every equivariant finitary representation }\Phi(Y)\overset{\mathrm{law}}=X,
 \quad \mathbb E[R_o(Y)]=\infty
\]
for every coding radius of that representation. The quantification concerns all representations of the law of \(X\), not merely a particular construction. As written, the question also permits a process having no finitary i.i.d. representation at all; the supplied record does not explicitly impose existence of such a representation.

The source does not determine whether the intended index set is \(\mathbb Z^d\), whether \(d\) is fixed, whether general transitive amenable graphs are allowed, or whether output states must be finite. The general graph framework above is therefore explanatory, not an attribution of these choices to the source.

\textbf{Supplied background and qualifications.} The dated curation states that every finitely dependent invariant process on a transitive amenable graph is a finitary factor of an i.i.d. process, that finite-mean or exponential coding-tail examples exist, and that the question remains unresolved. These are supplied literature assessments, not independently verified assertions in this draft.

The original excerpt then changes subject to Markov random fields on \(\mathbb Z^d\), describing their conditional distributions on finite sets as depending on the exterior configuration only through the boundary. It reports that a finite-state Markov chain is a finitary factor of an i.i.d. process exactly when it is ergodic, with ergodic explained there as aperiodic and irreducible. It also reports a countable-state criterion of exponential ergodicity, explained through exponential tails for return times to a fixed state. For the Ising model on \(\mathbb Z^d\), it attributes to van den Berg and Steif the assertion that the plus-state Gibbs measure is a finitary factor of an i.i.d. process exactly when the plus and minus states coincide, and identifies this with temperature at or above criticality. These neighboring claims have incompletely specified hypotheses and are not assumptions of the question. The excerpt ends during a further discussion of two ergodic Markov random fields with the same conditional distributions; no conclusion can be reconstructed from that unfinished sentence.
\subsection{Short English statement}
The question asks whether a finitely dependent invariant process can fail to have any finitary i.i.d. representation whose coding radius has finite expectation. This must fail for every possible representation, not just for one chosen coding. The supplied record leaves the index graph, invariance group, state spaces, and their quantifiers unspecified, so the displayed formulation remains conditional on those choices. The supplied open status requires independent review.
\subsection{Sources}
\begin{itemize}
\item Supplied local UnsolvedMath registry record, permanent ID \(30005024\), title \emph{Finite-Mean Coding Radius for Finitely Dependent Processes}; exact source problem code \texttt{catalogue entry}; collection: Oberwolfach Reports Open Problems. No proposer is named. The supplied proposed year is 2022.
\item Supplied report citation: \emph{Mini-Workshop: Descriptive Combinatorics, LOCAL Algorithms and Random Processes (2023)}, Oberwolfach Reports, pp.\ 429--455, DOI \texttt{10.4171/owr/2022/8}. Supplied source URL: \url{https://doi.org/10.4171/owr/2022/8}. The record describes these reports as documenting mathematical workshops and open-problem discussions.
\item Supplied literature source labeled ``Holroyd--Hutchcroft--Levy, finite-mean examples'': \url{https://doi.org/10.1214/19-AOP1363}. The supplied record gives no further bibliographic details or precise theorem statement.
\item Supplied literature source labeled ``OWR finitary-process report'': \url{https://doi.org/10.4171/owr/2022/8}.
\item Literature assessment checked at \texttt{2026-08-22}: the supplied curation labels the question open and reports the finitary-factor and coding-tail background described above. This is supplied dated curation derived from the OWR record, not an independently verified literature assessment; its status and cited claims require independent verification.
\item The registry labels the statement \texttt{corrected\_verified} and says it was reconstructed and checked against the official report. This records supplied provenance and does not constitute verification by this draft. The neighboring excerpt uses reference labels [6], [1], and [2], with [2] attributed to van den Berg and Steif, but supplies no corresponding bibliography entries.
\end{itemize}
\end{document}
